\color{unimain}{\subsection{Автоматическое повторное включение}}
\color{black}

\begin{enumerate}

\item
Автоматическое повторное включение (АПВ) предназначено для быстрого восстановления первоначального состояния электрической сети после аварийного отключения путём автоматического включения выключателя.

\item
Уставки АПВ представлены в таблице \ref{app_set:apv}.

\item
Выбор действующих режимов АПВ реализуется тремя входными сигналами ФСУ:
\begin{itemize}
    \item \textit{«ОНЭ1 + ННЭ2»} (включение с контролем отсутствия напряжения на линии и наличия напряжения на шинах) -- АПВ линии;
    \item \textit{«ОНЭ2 + ННЭ1»} (включение с контролем отсутствия напряжения на шинах и наличия напряжения на линии) -- АПВ энергообъекта Э1 (шин);
    \item \textit{«Без контролей»} (включение без контроля напряжений).
\end{itemize}

\item
Возможные режимы АПВ приведены в таблице \ref{desc_apv:tbl1}.

\small
\setlength{\extrarowheight}{0.05cm} % высота ячеек в таблице
\begin{longtable}{|>{\centering\arraybackslash}m{4.9cm}|>{\centering\arraybackslash}m{3.7cm}|>{\centering\arraybackslash}m{3.7cm}|>{\centering\arraybackslash}m{3.7cm}|}
  \caption{Выбор режима АПВ\hfill\vspace{-0.5\baselineskip}}\label{desc_apv:tbl1}\\
  \hhline{|-|-|-|-|}
  \rowcolor{gray!30}
  & \multicolumn{3}{c|}{Входные сигналы ФСУ} \\
  \hhline{|-|-|-|-|}
  \rowcolor{gray!30}
  \textit{Режим АПВ} & \textit{ОНЭ1+ННЭ2} & \textit{ОНЭ2+ННЭ1} & \textit{Без контролей} \\
  \hhline{|-|-|-|-|}

\endfirsthead

\caption*{\hspace{3pt}\emph{Продолжение таблицы \ref{desc_apv:tbl1}\hfill\vspace{-0.5\baselineskip}}} \\ % сделано по ГОСТ 2.105 п.6.8.7
\hhline{|-|-|-|-|}
\rowcolor{gray!30}
  & \multicolumn{3}{c|}{Входные сигналы ФСУ} \\
\hhline{|-|-|-|-|}
\endhead
\hhline{|-|-|-|-|}

\endfoot
\endlastfoot

КС(УС) 			  			            & 0 & 0 & 0 \\  \hline
ОНЭ1+ННЭ2 / КС(УС) 			        & 1 & 0 & 0 \\  \hline
ОНЭ2+ННЭ1 / КС(УС) 			        & 0 & 1 & 0 \\  \hline
ОНЭ1+ННЭ2 / ОНЭ2+ННЭ1 /  КС(УС) & 1 & 1 & 0 \\  \hline
Без контролей 				          & 0/1 & 0/1 & 1 \\  \hline

\end{longtable}
\setlength{\extrarowheight}{0.0cm} % возврат параметра высота ячеек в таблице
\normalsize

\item
Как видно из таблицы \ref{desc_apv:tbl1}, режим АПВ с КС(УС) введен всегда и выводится посредством входного сигнала ФСУ \textit{«Без контролей»}. А режимы ОНЭ1+ННЭ2 и ОНЭ2+ННЭ1 могут быть введены одновременно. Режим без контролей применяется только для АПВ линии. Условия включения в различных режимах более подробно описаны в пункте \ref{func:ksn} настоящего руководства по эксплуатации.

\item
Пуск АПВ энергообъекта Э1 (линии) осуществляется по цепи несоответствия при аварийном отключении выключателя (по наличию входного сигнала ФСУ \textit{«КСВ/ РФ: Сигнал несоотв.»}). Срабатывание происходит при набранной выдержке времени готовности АПВ с контролем заданного режима и выполнения условий включения. Выдержка времени готовности АПВ после оперативного включения и первого цикла регулируется уставкой \textbf{«Тгот АПВ»}, после второго цикла -- \textbf{«Твосст гот. АПВ»}.

\item
При приеме сигнала о срабатывании ДЗШ (входной сигнал ФСУ \textit{«АПВ/ АПВ Э1: Сброс АПВ от ДЗШ»}) готовность АПВ линии сбрасывается.

\item
АПВ линии может быть выполнено двукратным, ввод второго цикла осуществляется программным ключом \textbf{«2-ой цикл Э1»}. Первый цикл осуществляется при выполнении условий включения с выдержкой времени \textbf{«Тср АПВл Э1 1ц»}. В случае сброса сигнала пуска АПВ при отсчете времени \textbf{«Тср АПВл Э1 1ц»} цепь логики сбрасывается, и схема АПВ возвращается в исходное состояние (готовность АПВ с набранной выдержкой готовности). Если первый цикл АПВ был успешным, то начинается набор выдержки времени готовности к повторному действию.

\item
Если после первого цикла АПВ за время \textbf{«Тгот АПВ»} (выдержка времени к повторному действию не успела набраться) повторно поступает сигнал пуска АПВ, то осуществляется второй цикл с выдержкой времени \textbf{«Тср АПВ Э1 2ц»}. В случае сброса сигнала пуска АПВ при отсчете времени \textbf{«Тср АПВ Э1 2ц»} цепь логики сбрасывается, и схема возвращается в режим набора выдержки времени готовности к повторному действию. Если второй цикл АПВ был успешным, то начинается набор выдержки времени готовности к повторному действию.

\item
Второй цикл АПВ может быть выведен от внешнего сигнала \textit{«АПВ/ АПВ Э1: Вывод АПВ 2ц»}.

\item
Для схем с двумя выключателями на присоединение предусмотрен входной сигнал ФСУ \textit{«АПВ/ АПВ Э1: Ведомый»} для оперативного изменения порядка включения выключателя от АПВ линии – \textbf{«ведущий»} или \textbf{«ведомый»}. По умолчанию выключатель принимается «ведущий». При наличии сигнала \textit{«АПВ/ АПВ Э1: Ведомый»} время включения от АПВ линии дополнительно увеличивается на выдержку времени \textbf{«Тср АПВ Э1 доп.»}.

\item
Длительность импульса включения выключателя от АПВ линии задается уставкой \textbf{«Тимп АПВ Э1»}.

\item
Алгоритм работы АПВ энергообъекта Э1 (линии) приведен на рисунке \ref{pict_apv:APV_LAP}. 

\medskip
\begin{figure}[H]
\centering
\includegraphics[width=0.9\textwidth,height=0.9\textheight,keepaspectratio]{\fbpath/10.01 АУВ/АПВ/img/APV_LAP.pdf}
\captionof{figure}{Функциональная схема логики АПВ Э1}\label{pict_apv:APV_LAP}
\end{figure}

\item
Пуск АПВ энергообъекта Э2 (шин) осуществляется также по цепи несоответствия при аварийном отключении выключателя (по наличию логического сигнала \textit{«КСВ/ РФ: Сигнал несоотв.»}). АПВ шин выполнено однократным и осуществляется при выполнении условий включения с выдержкой времени \textbf{«Тср АПВ Э2»}. Длительность импульса включения выключателя от АПВ шин задается уставкой \textbf{«Тимп АПВ Э2»}. При приеме сигнала о срабатывании защит (входной сигнал ФСУ \textit{«АПВ/ АПВ Э2: Сброс АПВ от ср. защит»}) готовность АПВ шин сбрасывается.

\item
Алгоритм работы АПВ энергообъекта Э2 (шин) приведен на рисунке \ref{pict_apv:APV_CSH}.

\medskip
\begin{figure}[H]
\centering
\includegraphics[width=0.9\textwidth,height=0.9\textheight,keepaspectratio]{\fbpath/10.01 АУВ/АПВ/img/APV_CSH.pdf}
\captionof{figure}{Функциональная схема АПВ Э2}\label{pict_apv:APV_CSH}
\end{figure}


\end{enumerate}